\section{Runtime Analysis}
\label{sec:runtime}

\subsection{Theoretical analysis}

Each bisection step of METIS recursive bisection on a planar graph
with $|E_i|$ edges and $|V_i|$ vertices involves:

\begin{enumerate}
  \item Building the CSR graph representation: $O(|E_i|)$ time
  \item Heavy-edge matching for coarsening: $O(|E_i|)$ time per level,
    $O(\log |V_i|)$ coarsening levels
  \item Greedy initial partitioning on the coarsened graph: $O(|V_i|/c)$
    for coarsening ratio $c \approx 20$
  \item Fiduccia--Mattheyses refinement~\citep{fiduccia1982linear}:
    $O(|E_i| \cdot \texttt{niter})$ per level
\end{enumerate}

For a planar graph, $|E_i| = O(|V_i|)$ at every level.
The bisection tree has $\lceil \log_2 k \rceil$ levels.
At level $\ell$, there are $2^\ell$ subproblems each of expected size
$O(n / 2^\ell)$.  Total work at level $\ell$:

\[
  2^\ell \cdot O\!\left(\frac{n}{2^\ell} \cdot \log\!\left(\frac{n}{2^\ell}\right)\right)
  = O(n \log n).
\]

Summing over $\lceil \log_2 k \rceil$ levels:

\[
  T_{\mathrm{RB}}(n, k) = O(n \log n \cdot \log k).
\]

In the experiments, $\texttt{niter} = 100$ (a fixed constant).
For $n \leq 8{,}057$ (California), $\log n \approx 13$, so the FM
refinement runs for roughly $100/13 \approx 7.7$ times as many iterations
as the $O(\log n)$ asymptotic model.
Accordingly, the tightest correct statement of the runtime theorem uses
$\texttt{niter}$ explicitly as a constant multiplier.

\begin{theorem}[Runtime bound]
\label{thm:runtime}
METIS recursive bisection for the balanced $k$-cut redistricting problem
on a planar graph with $n$ vertices runs in
$O(\texttt{niter} \cdot n \log k)$ time with $O(n \log k)$ space,
where $\texttt{niter}$ is the number of FM refinement passes.
With the experimental setting $\texttt{niter} = 100$, this gives
$O(n \log k)$ up to a constant of 100.
\end{theorem}

\subsection{Empirical scaling across 50 states}

Table~\ref{tab:runtime} reports measured runtimes for representative states
from the 2020 Census.  The range for multi-district states spans from
Iowa (1,542 tracts, 4 districts) to California (8,057 tracts, 52 districts),
covering a $5\times$ range in $n$ among the states shown.%
\footnote{Single-district states ($k = 1$) are excluded
from the scaling analysis.  The seven excluded states are:
Alaska, Delaware, Montana, North Dakota, South Dakota, Vermont, and Wyoming.
For $k = 1$ there is no bisection step; the algorithm performs only
adjacency graph construction and returns immediately.
Runtime for these states is approximately 2~ms regardless of~$n$,
consistent with $O(n)$ graph-building time and a negligible constant.
Including them in the $T \sim n^b$ regression would bias the exponent
estimate by anchoring low-$n$ observations at near-zero runtimes.
The remaining 43 states ($k \geq 2$) form the OLS sample.}

\begin{table}[htb]
\centering
\caption{Empirical runtime scaling on 2020 Census tract data
  (mean of 10 runs $\pm$ std.\ dev., single-threaded, AMD Ryzen 9 5950X).
  Single-district states ($k = 1$) excluded; see text.}
\label{tab:runtime}
\begin{tabular}{lrrrc}
\toprule
State & Tracts ($n$) & Districts ($k$) & Runtime (ms) & ms/tract \\
\midrule
Iowa       & 1,542 &  4 & $18.4 \pm 0.6$ & 0.012 \\
Oregon     & 2,049 &  6 & $28.7 \pm 0.9$ & 0.014 \\
Wisconsin  & 1,921 &  8 & $31.2 \pm 1.1$ & 0.016 \\
Virginia   & 2,748 & 11 & $47.8 \pm 1.5$ & 0.017 \\
Michigan   & 2,813 & 13 & $51.4 \pm 1.7$ & 0.018 \\
Georgia    & 3,053 & 14 & $58.6 \pm 2.0$ & 0.019 \\
Ohio       & 2,884 & 15 & $53.2 \pm 1.8$ & 0.018 \\
Florida    & 4,245 & 28 & $98.7 \pm 3.2$ & 0.023 \\
New York   & 4,919 & 26 & $121.3 \pm 4.1$ & 0.025 \\
Texas      & 5,265 & 38 & $147.1 \pm 4.8$ & 0.028 \\
California & 8,057 & 52 & $198.4 \pm 5.9$ & 0.025 \\
\bottomrule
\end{tabular}
\end{table}

Fitting $\log T = \log a + b \log n$ by OLS on the $k > 1$ subset
(43~states, 2020~Census) gives $b = 1.07 \pm 0.03$ ($R^2 = 0.984$),
confirming near-linear scaling.
The residual super-linearity (0.07 above linear) is attributable to the
$\log k$ factor: larger states tend to have more districts, so both $n$
and $k$ grow together.

\paragraph{Log \boldmath$k$ covariate test.}
To test whether the $\log k$ term in the theoretical bound
$T = O(n \log k)$ accounts for the residual super-linearity,
we extend the regression to include $\log(\log k)$ as a second covariate:
\[
  \log T \;=\; \log a \;+\; b\,\log n \;+\; c\,\log(\log k).
\]
Fitting this two-variable model by OLS on the same 43-state, 2020-Census
sample gives $b = 1.05 \pm 0.03$, $c = 0.08 \pm 0.15$, $R^2 = 0.985$.
The coefficient on $\log(\log k)$ is not significantly different from zero
(95\% CI: $[-0.07,\,0.23]$; $p = 0.41$).
Adding $\log(\log k)$ as a covariate does not significantly improve fit
($\Delta$-AIC $< 2$, $c = 0.08 \pm 0.15$), confirming that $n$ is the
dominant scaling factor.
The seven $k = 1$ states excluded from this analysis are Alaska,
Montana, North Dakota, South Dakota, Vermont, Delaware, and Wyoming,
as noted in the footnote above.
The log $k$ range among the 43 included states is narrow
($\log k \in [0.69,\,3.95]$, i.e.\ $k$ from 2 to 52), which limits the
power to detect a $\log k$ effect if it exists --- but the point estimate
$c \approx 0.08$ is small relative to $b \approx 1.05$, and the CI
comfortably includes zero.
We conclude that $n$ alone is the sufficient predictor of runtime at
census-tract scale.

\subsection{Cross-census scaling}

Table~\ref{tab:census} compares scaling exponents across census years,
fitted on $k > 1$ states only.
The exponent is stable across years despite changes in tract count and
district assignments.

\begin{table}[htb]
\centering
\caption{Empirical scaling exponent $b$ in $T = a \cdot n^b$ by census year
  ($k > 1$ states only; single-district states excluded)}
\label{tab:census}
\begin{tabular}{lrrr}
\toprule
Census year & $n$ (total tracts, $k{>}1$ states) & Scaling exponent $b$ & $R^2$ \\
\midrule
2000 & 65,443 & $1.08 \pm 0.04$ & 0.981 \\
2010 & 73,056 & $1.06 \pm 0.03$ & 0.983 \\
2020 & 74,134 & $1.07 \pm 0.03$ & 0.984 \\
\bottomrule
\end{tabular}
\end{table}

The consistency across census years ($b \approx 1.07$ in all three) gives
confidence that the algorithm will scale similarly for the 2030 Census,
which is expected to enumerate approximately 78,000--82,000 tracts.

\subsection{Space complexity}

The bisection tree requires $O(n \log k)$ space to store the partition
hierarchy, as each of the $n$ tracts appears in one subproblem at each
of the $\lceil \log_2 k \rceil$ levels.

Peak memory during coarsening is $O(n)$ per subproblem (the coarsened
graph never exceeds the original size).  Since subproblems at the same
level are processed sequentially, total peak memory is:

\[
  M_{\mathrm{peak}} = O(n \log k).
\]

\paragraph{Numerical example.}
For California ($n = 8{,}057$, $k = 52$), $\log_2 52 \approx 5.7$, giving
a peak memory requirement of approximately $8{,}057 \times 6 \approx 48{,}000$
node records.  At 64 bytes per node record (population, adjacency pointers,
partition label), this is under 3\,MB --- trivial on modern hardware.

For block-group resolution ($n \approx 239{,}000$ nationally), peak memory
per state is under 20\,MB, well within the 16\,GB available on the
reference workstation.
